\chapter{Analyse et spécification des besoins}
\label{ch:analyse-besoins}

\section*{Introduction du chapitre}
Ce chapitre identifie les acteurs du système, recense les besoins
fonctionnels et non fonctionnels, puis les formalise à l'aide du langage
UML.

\section{Identification des acteurs}

Le système distingue cinq rôles internes, organisés en une hiérarchie de
sécurité (\texttt{ROLE\_SUPER\_ADMIN > ROLE\_ADMIN}), à laquelle s'ajoutent
deux acteurs externes ne possédant pas de compte : le tuteur de stage
(accès par lien à usage unique) et le visiteur public (formulaire de
contact, infolettre, vérification de certificat).

\begin{table}[H]
  \centering
  \caption{Rôles applicatifs et responsabilités}
  \label{tab:roles}
  \begin{tabular}{@{}p{2.8cm}p{9.5cm}@{}}
    \toprule
    Rôle & Responsabilités \\
    \midrule
    \texttt{SUPER\_ADMIN} & Paramètres système globaux, en plus de tous les droits Admin. \\
    \texttt{ADMIN} (Directeur) & Gestion pédagogique complète, validation des UF sensibles, gestion des groupes et des campagnes email, en plus de tous les droits Formateur. \\
    \texttt{FORMATEUR} & Studio de cours, gestion des quiz, correction manuelle, messagerie, bourse à l'emploi. \\
    \texttt{ETUDIANT} & Espace apprenant : contenu, quiz, stage, certificat, messagerie. \\
    \texttt{SUPPORT} & Accès limité, essentiellement en lecture. \\
    \bottomrule
  \end{tabular}
\end{table}
\tabdesc{Ce tableau met en correspondance chaque rôle applicatif avec les
acteurs métier identifiés au Tableau~\ref{tab:acteurs-academie} du chapitre
précédent, et sert de base à la hiérarchie \texttt{@PreAuthorize} détaillée
au chapitre~\ref{ch:conception}.}

\section{Spécification des besoins fonctionnels}
\label{sec:besoins-fonctionnels}

Le système doit permettre de~:
\begin{enumerate}[label=\textbf{BF\arabic*.}]
  \item représenter fidèlement la structure pédagogique réelle du cursus
    (Formation, Année, Unité de Formation, Module, Leçon, Bloc de contenu)~;
  \item gérer un système d'évaluation complet~: plusieurs types de quiz
    selon leur portée (section, module, fin d'UF, fin d'année), plusieurs
    types de questions, avec correction automatique et correction
    manuelle convergeant vers un score final cohérent~;
  \item garantir un déblocage séquentiel du contenu pédagogique, vérifié
    côté serveur~;
  \item fournir à l'équipe pédagogique un espace d'administration complet
    (contenu, correction, suivi des apprenants, validations manuelles)~;
  \item automatiser la délivrance de certificats numériquement
    vérifiables et outiller le suivi du stage jusqu'à sa signature par un
    tuteur externe~;
  \item sécuriser transversalement la plateforme (authentification,
    séparation des rôles, protection contre les abus) et permettre un
    déploiement reproductible.
\end{enumerate}

Ces six besoins reprennent directement les six objectifs du projet
énoncés au chapitre~\ref{ch:contexte} ; ils sont repris sous forme de
cas d'utilisation en section~\ref{sec:modelisation-besoins}.

\section{Spécification des besoins non fonctionnels}

\begin{table}[H]
\centering
\begin{tabular}{@{}p{3.5cm}p{10.5cm}@{}}
\toprule
\textbf{Besoin} & \textbf{Exigence associée} \\
\midrule
Sécurité & Authentification JWT signée HS256 transportée par cookie
  \texttt{httpOnly} (protection XSS), 2FA optionnelle, mots de passe
  hachés (BCrypt), RBAC déclaratif complété par un contrôle de propriété
  au niveau service (voir \S\ref{sec:acces-proprietaire}). \\
Protection contre les abus & Limitation de débit Redis sur les endpoints
  sensibles (ex. 20 requêtes/5 minutes pour la génération de questions
  par IA), voir \S\ref{sec:rate-limiting}. \\
Fiabilité fonctionnelle & Les règles de progression et de contrôle
  d'accès sont vérifiées côté serveur, jamais uniquement suggérées côté
  interface. \\
Scalabilité & Backend \texttt{STATELESS} (aucune session en mémoire) ---
  l'état d'authentification est reconstruit à chaque requête à partir du
  seul jeton JWT, ce qui simplifie la mise à l'échelle horizontale. \\
Maintenabilité & Architecture en couches stricte (\texttt{Controller
  $\rightarrow$ Service $\rightarrow$ Repository $\rightarrow$ Entity}),
  migrations de schéma versionnées (Flyway), suite de 138 tests
  automatisés. \\
Portabilité & Backend et frontend conteneurisés (Docker), configuration
  exclusivement par variables d'environnement (principe \emph{Twelve-Factor
  Apps}). \\
\bottomrule
\end{tabular}
\caption{Besoins non fonctionnels}
\label{tab:bnf}
\end{table}
\tabdesc{Ces exigences transversales conditionnent les choix
d'architecture présentés au chapitre~\ref{ch:conception}. Elles servent
également de critères d'acceptation lors de la phase de tests
(chapitre~\ref{ch:tests-deploiement-bilan}).}

\section{Modélisation des besoins}
\label{sec:modelisation-besoins}

\subsection{Diagramme de cas d'utilisation}

La Figure~\ref{fig:cas-utilisation} synthétise l'ensemble des cas
d'utilisation identifiés à partir des 29 contrôleurs REST exposés par le
backend, organisés en quatre paquetages fonctionnels (Sécurité \& Accès,
Parcours \& Évaluations, Stage \& Diplomation, Administration \&
Services), avec la hiérarchie de généralisation des acteurs et les
relations \textit{include}/\textit{extend} entre cas d'utilisation.

\begin{figure}[H]
  \centering
  \includegraphics[width=\textwidth,height=0.85\textheight,keepaspectratio]{diagramme_cas_utilisation.png}
  \caption{Diagramme de cas d'utilisation}
  \label{fig:cas-utilisation}
\end{figure}

\section*{Conclusion du chapitre}
Les besoins ayant été recensés et formalisés, le chapitre suivant présente
les choix de conception --- architecture, modèle de données, sécurité ---
permettant d'y répondre.
